% GENERATED FILE. Edit canonical lessons, then run npm run generate:books.
% canonical-source-hash: fnv1a64:7cc6c692b8058e11
% canonical-lessons: JA-C290-utsuru, JA-C290-shinkon, JA-C290-furusato, JA-C290-miokuru, JA-C290-atarashii

\chapter{A New Home}
\label{ch:ja-a-new-home}
\hlchaptermodality{290}

\begin{chapteropening}
\textbf{By the end of this chapter:} \emph{I can say where a family moved and who saw them off.}

Complete the last lesson of chapter 290: I can say where a family moved and who saw them off.
\end{chapteropening}

\section[\texorpdfstring{\ja{うつる} (utsuru)}{うつる (utsuru)}]{\ja{うつる} (utsuru) — to move (house); to move over}

\label{lesson:JA-C290-utsuru}



\begin{quote}
Before the new one: say the Japanese for to invite formally, then the Japanese for a married couple.
\end{quote}

\subsection*{You'll want to know: \ja{うつる}}
\textbf{\ja{うつる}} — \emph{utsuru} — \textquotedblleft{}to move (house); to move over\textquotedblright{}.

\begin{quote}
\textbf{\ja{かぞくで} \ja{いなかに} \ja{うつりました。}} — \emph{kazoku de inaka ni utsurimashita} — \textquotedblleft{}Our family moved to the country.\textquotedblright{}
\end{quote}

Polite \textbf{\ja{うつります}}, \emph{utsurimasu}; a godan verb. Past (polite): \textbf{\ja{あにが} \ja{となりの} \ja{まちに} \ja{うつりました。}}, \emph{ani ga tonari no machi ni utsurimashita}.

\begin{cousinweb}
\textbf{\ja{うつる}} (\emph{utsuru}) is a native Japanese word (wago), inherited rather than borrowed from Chinese; its kanji was attached later as a kun'yomi (JMdict).
\end{cousinweb}

\begin{grammarlens}[title={What you've built}]
One more word for telling what happened.
\end{grammarlens}

\subsection*{Your turn}
\begin{itemize}
\raggedright
  \item \emph{Say it:} \emph{utsuru}
  \item \emph{Say it:} \emph{utsuru}, once more
  \item \emph{Say it:} \emph{kazoku de inaka ni utsurimashita} — Our family moved to the country.
  \item \emph{Recall:} say the Japanese for to invite formally, then the Japanese for a married couple, then say \emph{utsuru} again
\end{itemize}

\begin{tcolorbox}[breakable,colback=teal!4,colframe=teal!35!black,title={Before you move on}]
What does \ja{うつる} mean? (\textquotedblleft{}To move (house); to move over\textquotedblright{}.) Say it once more.
\end{tcolorbox}
\section[\texorpdfstring{\ja{しんこん} (shinkon)}{しんこん (shinkon)}]{\ja{しんこん} (shinkon) — newly married}

\label{lesson:JA-C290-shinkon}



\begin{quote}
Before the new one: say the Japanese for a married couple, then the Japanese for to move (house).
\end{quote}

\subsection*{You'll want to know: \ja{しんこん}}
\textbf{\ja{しんこん}} — \emph{shinkon} — \textquotedblleft{}newly married\textquotedblright{}.

\begin{quote}
\textbf{\ja{ふたりは} \ja{しんこんです。}} — \emph{futari wa shinkon desu} — \textquotedblleft{}The two are newly married.\textquotedblright{}
\end{quote}

\begin{cousinweb}
\textbf{\ja{しんこん}} (\emph{shinkon}) is Sino-Japanese (kango): two Chinese characters read with their on'yomi, readings borrowed from Chinese. They are \emph{shin}, \textquotedblleft{}new\textquotedblright{} (Middle Chinese \emph{sin}) and \emph{kon}, \textquotedblleft{}marriage\textquotedblright{} (\emph{xwon}) (KANJIDIC2; Middle Chinese in Baxter's transcription).
\end{cousinweb}

\begin{grammarlens}[title={What you've built}]
One more word for telling what happened.
\end{grammarlens}

\subsection*{Your turn}
\begin{itemize}
\raggedright
  \item \emph{Say it:} \emph{shinkon}
  \item \emph{Say it:} \emph{shinkon}, once more
  \item \emph{Say it:} \emph{futari wa shinkon desu} — The two are newly married.
  \item \emph{Recall:} say the Japanese for a married couple, then the Japanese for to move (house), then say \emph{shinkon} again
\end{itemize}

\begin{tcolorbox}[breakable,colback=teal!4,colframe=teal!35!black,title={Before you move on}]
What does \ja{しんこん} mean? (\textquotedblleft{}Newly married\textquotedblright{}.) Say it once more.
\end{tcolorbox}
\section[\texorpdfstring{\ja{ふるさと} (furusato)}{ふるさと (furusato)}]{\ja{ふるさと} (furusato) — one's home town; one's birthplace}

\label{lesson:JA-C290-furusato}



\begin{quote}
Before the new one: say the Japanese for to move (house), then the Japanese for newly married.
\end{quote}

\subsection*{You'll want to know: \ja{ふるさと}}
\textbf{\ja{ふるさと}} — \emph{furusato} — \textquotedblleft{}one's home town; one's birthplace\textquotedblright{}.

\begin{quote}
\textbf{\ja{つまが} \ja{ふるさとに} \ja{かえる} \ja{ことを} \ja{言}\ja{い}\ja{出}\ja{しました。}} — \emph{tsuma ga furusato ni kaeru koto o iidashimashita} — \textquotedblleft{}My wife suggested going back to her home town.\textquotedblright{}
\end{quote}

\begin{cousinweb}
\textbf{\ja{ふるさと}} (\emph{furusato}) is a native Japanese word (wago), inherited rather than borrowed from Chinese, built from \emph{furu}, \textquotedblleft{}old\textquotedblright{} (from \emph{furui}), and \emph{sato}, \textquotedblleft{}village\textquotedblright{} (JMdict, KANJIDIC2).
\end{cousinweb}

\begin{grammarlens}[title={What you've built}]
One more word for telling what happened.
\end{grammarlens}

\subsection*{Your turn}
\begin{itemize}
\raggedright
  \item \emph{Say it:} \emph{furusato}
  \item \emph{Say it:} \emph{furusato}, once more
  \item \emph{Say it:} \emph{tsuma ga furusato ni kaeru koto o iidashimashita} — My wife suggested going back to her home town.
  \item \emph{Recall:} say the Japanese for to move (house), then the Japanese for newly married, then say \emph{furusato} again
\end{itemize}

\begin{tcolorbox}[breakable,colback=teal!4,colframe=teal!35!black,title={Before you move on}]
What does \ja{ふるさと} mean? (\textquotedblleft{}One's home town; one's birthplace\textquotedblright{}.) Say it once more.
\end{tcolorbox}
\section[\texorpdfstring{\ja{みおくる} (miokuru)}{みおくる (miokuru)}]{\ja{みおくる} (miokuru) — to see (someone) off}

\label{lesson:JA-C290-miokuru}



\begin{quote}
Before the new one: say the Japanese for newly married, then the Japanese for one's home town.
\end{quote}

\subsection*{You'll want to know: \ja{みおくる}}
\textbf{\ja{みおくる}} — \emph{miokuru} — \textquotedblleft{}to see (someone) off\textquotedblright{}.

\begin{quote}
\textbf{\ja{えきで} \ja{ともだちを} \ja{みおくりました。}} — \emph{eki de tomodachi o miokurimashita} — \textquotedblleft{}I saw my friend off at the station.\textquotedblright{}
\end{quote}

Polite \textbf{\ja{みおくります}}, \emph{miokurimasu}; a godan verb. Past (polite): \textbf{\ja{かぞくが} \ja{あにを} \ja{みおくりました。}}, \emph{kazoku ga ani o miokurimashita}.

\begin{cousinweb}
\textbf{\ja{みおくる}} (\emph{miokuru}) is a native Japanese word (wago), inherited rather than borrowed from Chinese, built from \emph{mi}, \textquotedblleft{}seeing\textquotedblright{} (from \emph{miru}), and \emph{okuru}, \textquotedblleft{}to send\textquotedblright{} (JMdict, KANJIDIC2).
\end{cousinweb}

\begin{grammarlens}[title={What you've built}]
One more word for telling what happened.
\end{grammarlens}

\subsection*{Your turn}
\begin{itemize}
\raggedright
  \item \emph{Say it:} \emph{miokuru}
  \item \emph{Say it:} \emph{miokuru}, once more
  \item \emph{Say it:} \emph{eki de tomodachi o miokurimashita} — I saw my friend off at the station.
  \item \emph{Recall:} say the Japanese for newly married, then the Japanese for one's home town, then say \emph{miokuru} again
\end{itemize}

\begin{tcolorbox}[breakable,colback=teal!4,colframe=teal!35!black,title={Before you move on}]
What does \ja{みおくる} mean? (\textquotedblleft{}To see (someone) off\textquotedblright{}.) Say it once more.
\end{tcolorbox}
\section[\texorpdfstring{\ja{あたらしい} (atarashii)}{あたらしい (atarashii)}]{\ja{あたらしい} (atarashii) — new}

\label{lesson:JA-C290-atarashii}



\begin{quote}
Before the new one: say the Japanese for one's home town, then the Japanese for to see (someone) off.
\end{quote}

\subsection*{You'll want to know: \ja{あたらしい}}
\textbf{\ja{あたらしい}} — \emph{atarashii} — \textquotedblleft{}new\textquotedblright{}.

\begin{quote}
\textbf{\ja{あたらしい} \ja{うちは} \ja{えきに} \ja{ちかいです。}} — \emph{atarashii uchi wa eki ni chikai desu} — \textquotedblleft{}The new house is close to the station.\textquotedblright{}
\end{quote}

An i-adjective.

\begin{cousinweb}
\textbf{\ja{あたらしい}} (\emph{atarashii}) is a native Japanese word (wago), inherited rather than borrowed from Chinese; its kanji was attached later as a kun'yomi (JMdict).
\end{cousinweb}

\begin{grammarlens}[title={What you've built}]
One more word for telling what happened.
\end{grammarlens}

\subsection*{Your turn}
\begin{itemize}
\raggedright
  \item \emph{Say it:} \emph{atarashii}
  \item \emph{Say it:} \emph{atarashii}, once more
  \item \emph{Say it:} \emph{atarashii uchi wa eki ni chikai desu} — The new house is close to the station.
  \item \emph{Recall:} say the Japanese for one's home town, then the Japanese for to see (someone) off, then say \emph{atarashii} again
\end{itemize}

\begin{tcolorbox}[breakable,colback=teal!4,colframe=teal!35!black,title={Before you move on}]
What does \ja{あたらしい} mean? (\textquotedblleft{}New\textquotedblright{}.) Say it once more.
\end{tcolorbox}
